% Fragmento público generado desde propietarios íntegros.
% Residencia: 52.9.2.
% Fuente: ../incoming/radion_monografia_20260827/manuscrito/sections/10_minkowski_hodge_lorentz.tex:101-186
\subsubsection{Lorentzian Realization of Helical Trajectories}

In a uniform magnetic field, a particle of charge \(q\), mass \(m\), and relativistic factor \(\gamma\) has
\begin{equation}
\omega_c=\frac{|q|B}{\gamma m},
\qquad
r_L=\frac{p_\perp}{|q|B},
\qquad
\Delta z=\frac{2\pi p_\parallel}{|q|B}.
\label{eq:lorentz-helix}
\end{equation}
The cyclotron phase closes while the parallel motion advances: the orbit is helical. The parallel with \eqref{eq:H9-helix} is structural and can be formalized by a transducer that preserves phase, orientation, and step. The quotient \(q_9\) is not identified with the spatial coordinate \(z\); rather, the same architecture of return plus advance is recognized.

The Lorentz force is not derived from the image of a spiral by similarity.
It is derived by completing \eqref{eq:lorentz-chain}. The image helps explain
why circular phase and longitudinal memory are compatible; the field
form determines the dynamics.

\subsubsection{Confluence between action, area, information, and gravity}

\paragraph{Cartan--Holst Realization of Holonomy}

The gravitational promotion is expressed via a holonomy morphism.
\begin{equation}
\operatorname{Hol}_{\mathcal A}
\bigl(\mathfrak R_{\mathrm{grav}}(\gamma)\bigr)
=\rho_C\bigl(\operatorname{Hol}_{\TPK}(\gamma)\bigr).
\label{eq:holonomy-gravity}
\end{equation}
The right-hand side contains route memory; \(\rho_C\) realizes it in a gravitational connection. Cartan torsion
\[
T^I=D_\omega e^I
\]
appears in the codomain. The global torsion \(\Delta_4\) is an upstream
coordinate that can feed the map, but it is not renamed as \(T^I\) without
\eqref{eq:holonomy-gravity}.

\paragraph{Transduction of the action area into informational area}

The radion fixes one quantum of action per phase. The Barbero hexad--octad
reader subsequently uses the channels \(q_\pm\), the incidence
\(90/120\), and the functional \eqref{eq:barbero-núcleo}. The causal direction is
\[
(A,C^*,\hret)\to q_\pm\to K_{6\to8}
\to\gamma_{\mathrm{BI}}\to\text{area/information}.
\]
The tritic alphabet provides the unit of entropy \(\log3\), and a cut of
\(81\) channels provides \(81\log3\). This entropy is not to be confused with
\(\log2\), with \(\log\varphi\), with a microcanonical entropy, or with a
von Neumann entropy unless the state is specified.

The physical synthesis is affirmative but typed: the radion supplies the oriented unit of action that a holonomy can transport; Barbero publishes the area normalization; the causal screen and Hodge make it possible to realize fields; Cartan--Holst receives the holonomy in gravitational geometry.

\begin{statusbox}
The scalar radion equation is not by itself an Einstein equation. Its
contribution to gravity is to provide an exact section of the architecture of
action, orientation, and holonomy consumed by the gravitational realization.
\end{statusbox}

\subsubsection{Wick Rotation as Analytic Reader Change}

In the conventional formulation, Wick rotation relates a Lorentzian time
variable to a Euclidean coordinate through complex continuation. In HMT,
the root \(J^2=-I\), the involutive map \(R^2=I\), and the threshold
\(N^2=0\) already exist. The change of reader
\[
J\longleftrightarrow R
\]
transports a common/oriented pair between elliptic publication and
hyperbolic publication. Complex continuation subsequently recognizes that
operation in analytic coordinates.

This explains why the circle, the action ellipse, and the hyperbolic interior can
belong to the same state without being the same metric. The circle uses \(J\) for
the phase; the shape of the ellipse uses a symplectic compression generated by \(R\); Klein measures
the interior by means of the cross-ratio; the Minkowski screen realizes causality
in another codomain.
